\documentclass[11pt]{article}
\usepackage{amsmath,amssymb}
\usepackage[margin=1in]{geometry}
\usepackage{booktabs}
\usepackage{tabularx}
\usepackage{hyperref}
\setlength{\emergencystretch}{2em}

\title{General Scientific Identity and Integration Superiority:\\
Building Globally Faithful Knowledge from Raw Heterogeneous Science}
\author{ORION Paper III Ultimate-Successor Working Manuscript}
\date{August 2026}

\begin{document}
\maketitle

\begin{abstract}
Scientific integration fails when lexical, schema, graph, or embedding similarity is mistaken for scientific identity. ORION Paper III targets the stronger problem that remains after modern scientific knowledge-graph construction, ontology matching, heterogeneous entity matching, provenance-aware governance, multi-source fusion, and KG-RAG are granted to the comparator: across heterogeneous scientific domains, source types, representations, temporal regimes, and downstream responsibilities, can ORION decide more faithfully what scientific objects are the same, different, plural, incompatible, or unresolved, while improving downstream science? The predecessor information-equivalent typed-product tie is preserved as evidence that representation centralization alone cannot establish superiority. P3-U therefore focuses on acquiring and applying the load-bearing scientific identity coordinates themselves: construct, measurement, context, causal role, temporal validity, provenance/stance, uncertainty, and representation sufficiency. Protected evaluation separates extraction error from identity-rule error, preserves genuine plurality and annotation ambiguity, requires source round-trip, and measures downstream utility under incomplete-knowledge controls. The intended terminal is \textbf{GENERAL\_SCIENTIFIC\_IDENTITY\_INTEGRATION\_SUPERIORITY}, including at least one failure-derived identity coordinate or integration mechanic that transfers to fresh scientific regimes.
\end{abstract}

\section{Largest scientific question}
\begin{quote}
Can ORION outperform the strongest donor-complete scientific integration systems at deciding what scientific objects are the same, different, plural, incompatible, or unresolved, and can that advantage improve downstream science across domains and representations?
\end{quote}

The baseline already includes strong structure. SciGraph-LLM extracts provenance-bound claims and evidence from scientific papers \cite{scigraphllm2026}. Oarga et al. generate scientific ontologies and KGs with open LLMs \cite{oarga2026}. OntoDup treats scholarly entity matching as a governed, evidence/provenance/version-aware decision process rather than a bare pairwise classifier \cite{ontodup2026}. GenOM enriches ontology concepts with generated descriptions and combines semantic, embedding, and lexical evidence for matching \cite{genom2026}. A 2026 survey and benchmark of heterogeneous entity matching explicitly separates representation and semantic heterogeneity and shows that test-time heterogeneity is a major stressor \cite{heterogeneity2026}. P3-U absorbs these capabilities and moves upward to the scientific identity responsibility relation itself.

\section{Upward-generalization principle}
Any donor that solves a P3-U subclass is included in the strongest comparator. ORION earns a broader contribution only through behavior such as:
\begin{enumerate}
    \item superior identity decisions after coordinate-equalized extraction;
    \item superior raw-to-identity performance across domains and source types;
    \item correct routing among \textsc{Glue}, \textsc{Distinct}, \textsc{Plural}, \textsc{Obstruction}, and \textsc{Unresolved};
    \item robustness to schema/representation changes with valid transport or reopening;
    \item higher downstream scientific synthesis/QA utility under incomplete-knowledge controls; and
    \item discovery of a new identity coordinate or integration mechanic that expands protected reach and transfers beyond the originating failure family.
\end{enumerate}
Generic KG extraction, ontology alignment, entity matching, provenance governance, and multi-source graph fusion are donor-owned.

\section{Immutable predecessor boundary}
The information-equivalent typed-product tie remains immutable. It rules out a claim that ORION wins merely because its representation is centralized, graph-shaped, or typed. The successor asks whether ORION can acquire the missing scientific identity relation under raw heterogeneous evidence, recognize when the available state is insufficient for that responsibility, and reopen exactly what is required.

\section{Identity relation and evidence coordinates}
P3-U evaluates
\[
I(x,y\mid c,e)\in\{\textsc{Glue},\textsc{Distinct},\textsc{Plural},\textsc{Obstruction},\textsc{Unresolved}\},
\]
where $c$ is scientific context and $e$ provenance-bound evidence.

Candidate coordinates include construct definition, operationalization, measurement procedure, scale/unit, population/context, intervention, causal role, temporal validity, provenance, source stance, modality, uncertainty, and correspondence maps. No fixed list is assumed universal. A major objective is to identify which coordinates are load-bearing in an unseen regime and, only when the registered basis is insufficient, propose a new coordinate prospectively.

\section{Prospective raw-evidence benchmark}
The corpus is frozen before protected outcomes, spans multiple scientific domains and source forms, and uses independent double annotation with adjudication. Genuine disagreement remains a first-class stratum rather than being erased by forced consensus.

\begin{table}[h]
\centering
\small
\begin{tabularx}{\textwidth}{@{}lXX@{}}
\toprule
Family & Surface trap & Identity obligation \\
\midrule
Construct drift & same term, changed construct & avoid false glue \\
Measurement mismatch & same latent construct, incompatible operationalization & preserve distinction/obstruction \\
Context dependence & same relation under different regimes & condition identity on context \\
Synonymy & different language, same scientific object & avoid false split \\
One-to-many mapping & one coarse object maps to several refined objects & preserve plurality \\
Contradictory evidence & sources disagree materially & calibrate obstruction/unresolved \\
Temporal change & identity relation changes after revision/discovery & bind validity in time \\
Provenance/stance & claim and criticized/negated claim share vocabulary & preserve stance \\
State insufficiency & compact state omits a decisive coordinate & refuse unsupported integration \\
Cross-representation transport & equivalent science appears through different schemas & require a valid bridge \\
\bottomrule
\end{tabularx}
\caption{Mandatory identity families. Every positive merge has a matched false-merge control and every distinction has a false-split/synonymy control.}
\end{table}

At least one held-out domain must require a load-bearing identity coordinate or mechanic absent from the development families. Otherwise mechanism-expansion claims are circular.

\section{Causal separation of extraction and identity}
A false merge/split is not automatically an identity-rule failure. Each protected case is evaluated in two complementary arms:
\begin{enumerate}
    \item \textbf{coordinate-equalized}: candidate systems receive the same protected extracted coordinates, isolating the identity rule;
    \item \textbf{raw end-to-end}: systems must extract, normalize, reason, and justify the identity relation from source material.
\end{enumerate}
A raw-text failure that disappears under coordinate equalization is attributed to extraction/normalization before the identity rule is changed.

\section{Representation and responsibility stress test}
The same scientific case is presented through raw source evidence, canonical extracted state, task-compiled state, and responsibility-carrying state. If two views are independently certified semantics-preserving and sufficient for the identity responsibility, the promoted decision should be invariant. If a compact state supports QA or prediction but omits an identity-critical coordinate, the correct result is state insufficiency/\textsc{Unresolved}; raw evidence must remain recoverable.

Every promoted \textsc{Glue}, \textsc{Distinct}, \textsc{Plural}, or \textsc{Obstruction} decision must round-trip to the source spans or records carrying its load-bearing coordinates. A generated ontology label, embedding score, graph edge, or confidence value cannot self-authorize scientific identity.

\section{Causal Identity Coordinate Discovery}
The ChatGPT-led negative-to-positive successor family is \textbf{Causal Identity Coordinate Discovery (CICD)}. CICD treats each retained false merge, false split, plurality loss, or obstruction miss as competing hypotheses over extraction, construct, measurement, context/temporal, provenance/stance, state sufficiency, transport, or annotation non-identifiability.

The programme freezes a discriminator that changes one suspected coordinate while holding others fixed. A new scientific coordinate is proposed only after the registered coordinate basis and strong donor mappings fail prospectively on the frozen family. Generic feature discovery, metric learning, ontology induction, or representation learning is donor-owned; the residual is the scientific responsibility rule governing when a coordinate is necessary, what prior identity assertions must reopen, and whether the new coordinate transfers.

\section{Donor-complete baselines}
All systems receive matched model, evidence access, candidate-visible information, and resource budgets. Required comparators include:
\begin{enumerate}
    \item lexical/embedding entity matching;
    \item heterogeneity-aware entity matching \cite{heterogeneity2026};
    \item GenOM-style ontology matching \cite{genom2026};
    \item provenance-aware scientific KG extraction \cite{scigraphllm2026};
    \item scientific ontology/KG generation \cite{oarga2026};
    \item OntoDup-style governed identity materialization \cite{ontodup2026};
    \item LLM raw-text pair judge and RAG/entity-canonicalization system;
    \item multi-source KG fusion and incomplete-KG reasoning controls;
    \item compiled-state/responsibility-aware integration;
    \item donor-complete product containing all preceding mechanisms and allowed to route among them;
    \item strongest runnable scientific integration system available at protocol freeze;
    \item oracle identity-coordinate ceiling for analysis only.
\end{enumerate}

\section{Primary hypotheses}
\paragraph{H1, generalized identity superiority.}
Across the heterogeneous benchmark, ORION reduces false scientific merges versus the strongest runnable donor-complete comparator by a prospectively frozen margin.

\paragraph{H2, non-compensatory plurality guard.}
The H1 gain is invalid if false splits or legitimate plurality loss exceed frozen non-inferiority guards.

\paragraph{H3, downstream scientific superiority.}
The ORION portrait yields higher protected synthesis/QA utility under matched downstream models and incomplete-knowledge/leakage controls.

\paragraph{H4, cross-domain and cross-representation transfer.}
The H1--H3 direction survives held-out domains, schemas, source forms, representations, and independent annotation/implementation.

\paragraph{H5, mechanism expansion.}
At least one retained negative/tied family yields a new identity coordinate or integration mechanic that expands protected reach beyond the pre-freeze donor-complete basis and transfers to fresh families.

\section{Statistics and annotation authority}
The scientific item/pair or prospectively declared dependent cluster, not repeated model samples or judges, is the independent unit. Before protected outcomes, freeze sampling, cluster structure, family weights, adjudication procedure, agreement statistic, superiority/non-inferiority margins, multiplicity handling, interval method, and treatment of unresolved cases.

False merge, false split, plurality loss, obstruction miss, and unresolved calibration are reported separately with uncertainty. A single F1 score cannot compensate false scientific glue with easier true matches. Expert disagreement strata remain visible after adjudication.

\section{Framework correspondence}
Current framework substrate is semantically aligned with this manuscript. `ScientificMeaningProjection.v1` records source span, predicate, referents, constructs, measurements, temporal context, attribution, discourse relation, polarity, modality, assumptions, and unresolved ambiguity. `compare_meaning` conservatively returns distinct referent/construct/measurement, contextual difference, contradiction, or unresolved rather than defaulting to equivalence. It is explicitly proposal-level: extraction does not create scientific authority.

`MeasurementRelation.rakl-v1`, `SimilarityWitness.rakl-v1`, `GeneratorTransport.rakl-v1`, `EpistemicContextCompiler.rakl-v1`, and existing P3 independent-evaluator infrastructure are supporting research substrate. Naturalistic identity-coordinate acquisition, CICD, and GENERAL\_SCIENTIFIC\_IDENTITY\_INTEGRATION\_SUPERIORITY remain prospective and are not represented as current runtime capability.

\section{Adversarial controls and negative recursion}
Mandatory attacks include same terminology/different construct, different terminology/same construct, schema compatibility with incompatible measurement semantics, identical graph structure with different provenance/stance, compact state sufficient for QA but insufficient for identity, source-round-trip failure, answer leakage in downstream QA, and a donor-complete dynamic router.

Every material negative is preserved and assigned a primary responsibility. The programme then freezes a discriminator, tests materially different successor mechanisms, and requires fresh-family transfer plus downstream benefit. Gold relabeling after viewing errors, post-hoc pair deletion, baseline weakening, and replacing asymmetric harm with a single aggregate score are prohibited.

ORION issue \#651 remains the authority-bearing P3-U lane. CICD is the first ChatGPT-led recovery family.

\section{Claim ladder and current status}
The intended progression is protocol frozen; protected identity superiority; downstream scientific superiority; cross-domain/cross-representation replication; and finally \textbf{GENERAL\_SCIENTIFIC\_IDENTITY\_INTEGRATION\_SUPERIORITY} with H5 mechanism expansion. This TeX is prospective and does not rewrite the predecessor tie or claim broad raw-literature superiority before protected evidence exists.

\begin{thebibliography}{12}
\bibitem{scigraphllm2026}
I. P. Malashin et al., ``SciGraph-LLM: Automatic Knowledge Graph Construction from Scientific Papers,'' WSDM Companion 2026, pp. 48--57, doi:10.1145/3779211.3793169.

\bibitem{oarga2026}
A. Oarga, M. Hart, A. M. Bran, M. Lederbauer, and P. Schwaller, ``Scientific knowledge graph and ontology generation using open large language models,'' \emph{Digital Discovery}, vol. 5, pp. 1269--1279, 2026, doi:10.1039/d5dd00275c.

\bibitem{ontodup2026}
J. Galan-Mena, M. Lopez-Nores, D. Pulla-Sanchez, L. F. Guerrero-Vasquez, and J. P. Salgado-Guerrero, ``OntoDup: Governance-Aware Entity Matching for Scholarly Knowledge Graph Deduplication,'' \emph{Information}, vol. 17, no. 4, 325, 2026, doi:10.3390/info17040325.

\bibitem{genom2026}
Y. Song, J. Chen, and R. A. Schmidt, ``GenOM: ontology matching with description generation and large language models,'' \emph{World Wide Web}, vol. 29, article 29, 2026, doi:10.1007/s11280-026-01413-y.

\bibitem{heterogeneity2026}
M. H. Moslemi, A. Mousavi, B. Behkamal, and M. Milani, ``Heterogeneity in entity matching: A survey and experimental analysis,'' \emph{Data \& Knowledge Engineering}, vol. 164, 102575, 2026, doi:10.1016/j.datak.2026.102575.
\end{thebibliography}

\end{document}
